\chapter{Scopes, and the Fixed Points that Generate Them}
\label{ch:scope}
% ===================================================================

\begin{quote}\itshape\small
The previous two chapters built a logic: a language of formulae generated
from the calculus rather than imposed on it, whose sentences denote
equivalence classes of behaviour. This chapter puts that language to a use
it has not yet been put to. A name predicate does not only describe a
name; it carves out a \emph{region} --- and a region is what everything
later in this book will need when it wants to say where something is, how
many of something there are, or what a learner owns. The machinery is
short. What it buys is that a region can be finitely described and
unboundedly large at the same time, which is the condition every
interesting application turns out to be in.
\end{quote}

\section{Regions, and why a list will not do}

In the rho calculus a channel is a name $\quo{P}$, and a name is a piece
of code at which something may reside. Location is therefore already a
computational notion here, in a way it is not in most formalisms: to say
where a process is, is to exhibit a program.

That makes the obvious definition of a region available at once. A region
is a finite set of channels, $\{c_1, \dots, c_k\}$, and anything one wants
to say about ``here'' is said by quantifying over that set. Counting is
counting members; ownership is membership; a boundary is the difference
between two such sets.

The definition works and it is too weak, for a reason that has nothing to
do with cardinality and everything to do with description. A finite set
must be given by listing it, and almost no region worth naming can be
listed. Consider three that this book will need:

\begin{itemize}\itemsep2pt
\item the channels belonging to a particular learner;
\item the channels of \emph{any} learner rooted at a given quotation;
\item the channels of an organism, and of its organs, and of their cells,
      and so on down.
\end{itemize}

\noindent
The first is finite but not known in advance; it is what a learner is
trying to establish about itself. The second is not finite. The third is
not finite and is not even flat --- it has levels, and the levels are the
point.

A learner can only work with a region it can \emph{state}, because
everything a learner traffics in is a formula and everything it does with
a formula is priced by that formula's structure. So the region wants to be
a formula too.

\section{Scope as a predicate}

\begin{definition}[Scope]
\label{sco:def:scope}
A \emph{scope} is a name predicate $\Nsp$. Its \emph{extension} is
\[
  \Ext(\Nsp) \;=\; \{\, n \;:\; n \models \Nsp \,\},
\]
a collection of channels, possibly infinite. We write $\quo{\phi}$ for the
scope satisfied by $\quo{P}$ exactly when $P \models \phi$, and read it
``names of $\phi$-things''.
\end{definition}

Nothing here is new machinery: $\quo{\phi}$ is the name predicate of the
namespace logic \cite{NamespaceLogic}, and the only move being made is to
take its extension seriously as a \emph{place}.

\begin{remark}[The extensional version is the flat case]
\label{sco:rmk:flat}
A finite set of channels $\{c_1,\dots,c_k\}$ is the scope
$\quo{\phi_1} \lor \dots \lor \quo{\phi_k}$, where $\phi_i$ is a
characteristic formula for $\drop{c_i}$. So nothing is given up in the
move to predicates. What is gained is that a scope may now be finitely
described and unboundedly large, and that a scope is the same kind of
object as a hypothesis --- which means the same economics applies to both.
\end{remark}

We will use \emph{namespace} for the object and \emph{scope} for the role
it plays when something is being located or counted in it. The two words
name the same thing seen from two sides, and both are already in use in
this book: the namespaces of \chComp are scopes in exactly this sense.

\section{Disjointness is structure}

A composite region ought to be harder to work in than its parts. It is
not, provided the parts are disjoint, and the reason is a small
proposition about parsing.

\begin{proposition}[Unique decomposition]
\label{sco:prop:parse}
Let $\quo{\phi}$ and $\quo{\psi}$ be scopes with
$\Ext(\quo{\phi}) \cap \Ext(\quo{\psi}) = \emptyset$, separated at grade
$0$ in the sense of \S\ref{com:S5}. Then every
$n \in \Ext(\quo{(\phi \mid \psi)})$ decomposes uniquely: there are $P, Q$
with $n = \quo{(P \mid Q)}$, $P \models \phi$, $Q \models \psi$, and $P,Q$
are determined up to structural congruence.
\end{proposition}

\begin{proof}[Proof sketch]
Existence is the reading of $\mid$ at formula level. For uniqueness,
suppose $n = \quo{(P \mid Q)} = \quo{(P' \mid Q')}$ with
$P, P' \models \phi$ and $Q, Q' \models \psi$. Quotation is injective, so
$P \mid Q \equiv P' \mid Q'$. Grade-$0$ separation says there is no
spanning redex between the parands, so by
Theorem~\ref{com:thm:factor} the free names of the $\phi$-part and the
$\psi$-part are disjoint. A structural congruence between two parallel
compositions whose parands have disjoint free names can only permute
parands within each side; hence $P \equiv P'$ and $Q \equiv Q'$.
\end{proof}

\begin{corollary}[Description is additive under disjointness]
\label{sco:cor:additive}
Under the hypotheses of Proposition~\ref{sco:prop:parse}, the description
length of the composite scope is the sum of those of its parts up to a
constant, and membership testing decomposes: a name is tested by splitting
it once and testing each part against its own predicate.
\end{corollary}

\begin{remark}[Overlap is where you pay, and it is the same payment as before]
\label{sco:rmk:overlap}
If the scopes overlap, the split is not unique, and a membership test must
search over splittings rather than perform one. The description of the
composite then fails to decompose, and the excess is exactly the defect of
the lax comparison map of \chComp: the interaction between two learners
and the ambiguity in parsing their joint namespace are one quantity, seen
once dynamically and once syntactically. This is the third appearance of
that pattern, after \chComp's $\lambda$ and \chEng's $\theta$, and its
recurrence is worth noticing --- in each case something that looks like a
failure of composition to be clean turns out to be the measure of what
composition achieved.
\end{remark}

So disjointness is not an obstacle that composition works around. It is
the certificate that makes a composite scope \emph{structured}, and
structure is what makes a large scope cheaper to work in than a small
unstructured one. This is the opposite of the usual intuition about
modularity, where separation is a discipline one pays for; here separation
is what one is buying.

\section{Fractality by fixed point}

Disjoint composition gives wider scopes. It does not by itself give deeper
ones --- scopes whose members are themselves composites of members, all
the way down. For that we want a fixed point, and the logic has the
apparatus already \cite{kozen}.

\begin{definition}[Generated scope]
\label{sco:def:gen}
Let $\phi, \psi$ be formulae and $X$ a scope variable. The
\emph{generated scope}
\[
  \Nsp \;=\; \mu X.\, \quo{\bigl( (\phi \lor X) \mid (\psi \lor X) \bigr)}
\]
is the least scope with
$\Ext(\Nsp) = \Ext\bigl(\quo{((\phi \lor \Nsp) \mid (\psi \lor \Nsp))}\bigr)$,
where a scope variable in a process position is read as ``the drop of a
name in $\Ext(X)$''.
\end{definition}

More elaborate generators are available and will be wanted --- one
predicate per skill rather than two, guarded recursion under several
variables, gradings carried through the unfolding --- and
\S\ref{scm:sec:composite} uses the four-predicate version. This one is
enough to make the points.

\begin{definition}[Stratum]
\label{sco:def:stratum}
The \emph{stratum} $\str(n)$ of a name $n \in \Ext(\Nsp)$ is the number of
unfoldings of $\mu X$ needed to derive $n \models \Nsp$. Atoms --- names of
$\phi$-things and $\psi$-things that are not themselves composites ---
have stratum $0$.
\end{definition}

\begin{figure}[t]
\centering
\begin{tikzpicture}[scale=0.92,
  atom/.style={circle,draw,minimum size=5mm,inner sep=0pt,font=\scriptsize},
  comp/.style={rectangle,draw,rounded corners=2pt,minimum height=5mm,
               inner sep=2.5pt,font=\scriptsize},
  lbl/.style={font=\scriptsize\itshape,black!60}]

\node[lbl,anchor=east] at (-0.4,0)   {stratum $0$};
\node[lbl,anchor=east] at (-0.4,1.5) {stratum $1$};
\node[lbl,anchor=east] at (-0.4,3.0) {stratum $2$};

\node[atom] (a) at (0.6,0) {$\phi$};
\node[atom] (b) at (2.0,0) {$\psi$};

\node[comp] (c1) at (0.3,1.5) {$\quo{(\phi\mid\psi)}$};
\node[comp] (c2) at (2.9,1.5) {$\quo{(\phi\mid\phi)}$};
\node[comp] (c3) at (5.4,1.5) {$\quo{(\psi\mid\psi)}$};

\node[comp] (d1) at (1.0,3.0) {$\quo{\bigl((\phi\mid\psi)\mid\psi\bigr)}$};
\node[comp] (d2) at (5.0,3.0) {$\quo{\bigl((\phi\mid\psi)\mid(\phi\mid\phi)\bigr)}$};

\draw[-{Stealth[length=1.6mm]},black!55] (a) -- (c1);
\draw[-{Stealth[length=1.6mm]},black!55] (b) -- (c1);
\draw[-{Stealth[length=1.6mm]},black!55] (a) -- (c2);
\draw[-{Stealth[length=1.6mm]},black!55] (b) -- (c3);
\draw[-{Stealth[length=1.6mm]},black!55] (c1) -- (d1);
\draw[-{Stealth[length=1.6mm]},black!55] (b) to[out=120,in=-30] (d1);
\draw[-{Stealth[length=1.6mm]},black!55] (c1) to[out=20,in=200] (d2);
\draw[-{Stealth[length=1.6mm]},black!55] (c2) -- (d2);

\draw[dashed,black!35] (-0.2,0.75) -- (7.6,0.75);
\draw[dashed,black!35] (-0.2,2.25) -- (7.6,2.25);

\node[lbl,anchor=west,align=left] at (7.8,1.5)
  {the generator is\\ the same size\\ at every depth};
\end{tikzpicture}
\caption{The extension of a generated scope, stratified. Every name at
stratum $j+1$ quotes a parallel composition of drops of names at stratum
$\le j$, and by Proposition~\ref{sco:prop:parse} that decomposition is
unique.}
\label{sco:fig:strata}
\end{figure}

\begin{table}[t]
\centering
\begin{tabular}{rrr}
\toprule
stratum $\le j$ & names & new at this stratum \\
\midrule
$0$ & $2$      & $2$ \\
$1$ & $5$      & $3$ \\
$2$ & $17$     & $12$ \\
$3$ & $155$    & $138$ \\
$4$ & $12{,}092$ & $11{,}937$ \\
$5$ & $73{,}114{,}280$ & $73{,}102{,}188$ \\
\bottomrule
\end{tabular}
\caption{Extension of the two-atom generator of
Definition~\ref{sco:def:gen}, counted up to structural congruence:
$u_0 = 2$ and $u_{j+1} = 2 + \binom{u_j+1}{2}$, since the composites at
stratum $\le j+1$ are unordered pairs with repetition drawn from the names
at stratum $\le j$.}
\label{sco:tab:growth}
\end{table}

The extension grows doubly exponentially while the generator does not grow
at all. That gap is the whole point of the construction, and
\S\ref{sco:sec:genlen} names the quantities on either side of it. First,
though, the property that makes generated scopes usable rather than merely
compact.

\begin{proposition}[Membership is decidable by descent]
\label{sco:prop:descent}
Let $\Nsp$ be a generated scope whose atomic predicates have decidable
satisfaction. Then for any name $n$, the question $n \models \Nsp$ is
decidable.
\end{proposition}

\begin{proof}[Proof sketch]
Test $n$ by unfolding: $n \models \Nsp$ iff $\drop{n}$ satisfies one of
the atoms, or $n = \quo{(P \mid Q)}$ with $\quo{P}$ and $\quo{Q}$ each
satisfying $\Nsp$. By Proposition~\ref{sco:prop:parse} the split is
unique, so no search is required. The recursive calls are on $\quo{P}$ and
$\quo{Q}$, and by the no-self-code theorem \cite{RHO} a name codes a
strictly smaller term, so $|P|, |Q| < |P \mid Q|$. The recursion is
therefore well-founded on term size and terminates, with depth bounded by
$\str(n)$.
\end{proof}

\begin{remark}[What is doing the work]
\label{sco:rmk:whatworks}
Two hypotheses carry this and neither is decoration. Disjointness makes
the split unique, so the descent is deterministic rather than a search;
without it the procedure branches over splittings and its cost is no
longer governed by term size. No-self-code makes the descent
well-founded, and that is a property of \emph{reflection} specifically ---
it is what prevents a name from being a member of a scope by way of
itself. So the computability here is not bought by restricting to finite
regions. It is bought by the calculus being reflective, which is the same
purchase that \S\ref{eng:S5.4} makes when it builds the tower of media and
\S\ref{sci:S12.2} makes when it puts the germ on a channel.
\end{remark}

\begin{remark}[$\mu$, not $\nu$, and this is a claim]
\label{sco:rmk:munu}
Definition~\ref{sco:def:gen} takes the least fixed point, and the choice
is not innocent. The greatest fixed point admits non-well-founded scopes
--- names whose membership is witnessed only by an infinite descent ---
and Proposition~\ref{sco:prop:descent} does not cover them, since the
argument from no-self-code bounds the descent by term size and a
non-well-founded witness has no such bound.

We take the view that this framework only ever wants $\mu$, and that the
reason is economic rather than conventional. A scope a learner can survey
is one whose membership test terminates; a test that does not terminate is
a cost with no receipt, and by \S\ref{sci:S4} an assay that cannot be
settled within budget is refuted rather than pending. A $\nu$-scope is
therefore a region nothing budgeted can knowingly inhabit. That is a
substantive claim about what kinds of places there are for a mind to be,
and we would rather state it than assume it; the alternative is recorded
as an open question, since a non-well-founded scope is exactly what one
would reach for to model an ecology with no bottom.
\end{remark}

\section{Generator length}
\label{sco:sec:genlen}

\begin{definition}[Generator length]
\label{sco:def:genlen}
The \emph{generator length} $\gen(\Nsp)$ of a generated scope is the
number of symbols in the $\mu$-formula defining it, with atomic predicates
counted by their own lengths.
\end{definition}

Table~\ref{sco:tab:growth} is the observation that $\gen$ and
$|\Ext(\Nsp)|$ are not the same kind of quantity and do not move together.
Chapter~\ref{sec:scope-manufacture} adds a third, the number of joining
steps it takes to \emph{build} an inhabitant of the scope, and shows that
the three stand in a chain. We flag the destination here because the
temptation, on meeting a compact description of an enormous structure, is
to conclude that the structure is cheap. It is not. The description is
cheap.

\section{What this is for}

Three uses, each in a later part, and each is a place where the book has
so far been managing with the extensional notion and straining.

\emph{Ownership.} \chEng defines an individual as a region whose internal
media are closed and whose boundary media are not, and observes that the
definition is a fixed-point condition rather than a construction. With
Definition~\ref{sco:def:gen} the circularity has a shape: what is being
solved for is a generator, and the levels of the resulting hierarchy are
its strata.

\emph{Counting.} Chapter~\ref{sec:assembly} counts copies of a process in
a namespace. That count needs a region, and
Remark~\ref{sco:rmk:flat} is why the region cannot in general be a list.
Chapter~\ref{sec:scope-manufacture} takes this up in earnest.

\emph{Knowing.} A learner's reach is bounded by what it can afford to
visit, and \S\ref{eng:S8} prices visiting. A scope is therefore not only a
description but a budget line, and the region that exists \emph{for a
given learner} is the part of the extension it can pay to survey. The
Prestige returns to this in \S\ref{sec:kant-affordable}, where the
distinction between what can be characterised and what can be afforded
turns out to be doing metaphysical work.

% ===================================================================
